% TÜBİTAK 2204-A Proje Raporu — DART-RIFT.
% `docs/rapor2204` içinden iki kez derle (çapraz başvurular için):
%     lualatex rapor.tex
% pdflatex DEĞİL: parçacık ölçekli şekiller pdflatex'in sabit ana belleğini
% tüketiyor. Şekiller `make` ile ayrıca üretilir.
\documentclass[11pt,a4paper]{article}
\usepackage{tubitak}

% Kimlik bilgisi varsayılan olarak GİZLİ (bkz. tubitak.sty). Kılavuz
% kimlikli sürüm isterse aşağıdaki satırın başındaki % kaldırılır.
% \kimlikliyazdirtrue

\projealani{Fizik \textbullet\ Astronomi ve Astrofizik}
\projekodu{TÜBİTAK 2204-A Lise Öğrencileri Araştırma Projeleri Yarışması}
\ekip{Yağız Ekrem Dalar \quad Feyzi Arda Salihoğlu \quad Nedim Mutlu Sezer}
\kurum{Pertevniyal Lisesi, Fatih/İstanbul}
\ustbilgi{Proje Raporu \textbar\ 12 Ekim 2026}

\title{Dimorphos'un İç Yapısının DART Momentum Aktarımından\\
       GPU Tabanlı SPH Benzetimleriyle Çıkarımı}
\date{12 Ekim 2026}

\hypersetup{
  pdftitle={Dimorphos'un İç Yapısının DART Momentum Aktarımından GPU
            Tabanlı SPH Benzetimleriyle Çıkarımı},
  pdfsubject={TÜBİTAK 2204-A proje raporu — Fizik / Astronomi ve Astrofizik}
}

\begin{document}
\maketitle

\begin{abstract}
NASA'nın DART görevi 2022'de Dimorphos asteroidine çarparak yörünge
dönemini ölçülebilir biçimde değiştirdi ve cisme aktarılan momentumu
$\beta$ katsayısıyla sınırladı. Bu tek sayıdan cismin \emph{iç yapısını}
geri okumak doğrudan bir ters problem değildir: yakın tarihli çalışmalar
birbirinden farklı iç yapıların aynı $\beta$ değerini ürettiğini
göstermiştir. Bu projede, çarpmayı GPU üzerinde düzgünleştirilmiş
parçacık hidrodinamiği (SPH) ile modelleyen bir motor yazıldı ve üç iç
yapı koordinatının --- blok gözenekliliği $\alpha_b$, matris bağdaşım
dayanımı $Y_0$, blok hacim kesri $f$ --- gözlemden ne kadarının
öğrenilebildiği sorusu, bir \emph{nokta tahmini} yerine
\emph{tanımlanabilirlik} sorusu olarak kuruldu. Sabit bir DART sahnesinde
geçerliliği denetlenmiş beş benzetim tamamlandı. Bağdaşım 10 Pa'dan
5000 Pa'ya çıkarıldığında $\beta$ 3,7480'den 1,8738'e düşmekte, fırlatılan
kütle 26,1 kat azalmaktadır. Denenen matris çekme yasasının açılması
fazla momentumu \%97,1 değiştirmekte; çarpma noktasının 25 m kaydırılması
ise $\beta$'yı yalnız \%2,2 değiştirirken fırlatı kütlesini \%18
değiştirmektedir. Ölçülen tüm kaldıraçlar tek ölçekte sıralandığında en
büyük etki, çıkarmaya çalıştığımız koordinatlardan değil, kararla
sabitlediğimiz bir model seçiminden gelmektedir. Üretim çıkarımının 96
noktalık tasarımı, olabilirlik tanımı ve beş karar kapısı veriye
bakılmadan önce yazılıp tarih damgasıyla kayda alınmıştır; bu rapor o
kayıtları ve pilot ölçümleri sunar, üretim posterioru ya da sayısal Hera
öngörüsü sunmaz.
\end{abstract}

\anahtarkelimeler{gezegen savunması, çarpma fiziği, düzgünleştirilmiş
parçacık hidrodinamiği, Bayesçi çıkarım, Dimorphos}

% =======================================================================
\section{Giriş}

\subsection{Problem durumu}
26 Eylül 2022'de NASA'nın DART (\emph{Double Asteroid Redirection Test})
uzay aracı, Didymos ikili asteroit sisteminin uydusu Dimorphos'a saniyede
yaklaşık 6,1 km hızla çarptı. Çarpma, Dimorphos'un Didymos çevresindeki
yörünge dönemini ölçülebilir biçimde kısalttı
\citep{daly2023}. Bu ölçüm, cisme aktarılan momentumun çarpıcının kendi
momentumunun kaç katı olduğunu veren \emph{momentum artırma katsayısı}
$\beta$'yı sınırladı \citep{cheng2023}. $\beta>1$ olması, fırlatılan
malzemenin geri tepmesinin çarpıcının kendi momentumuna eklendiği anlamına
gelir; yani $\beta$, doğrudan ne kadar malzemenin ne hızla atıldığına,
o da cismin \emph{iç yapısına} bağlıdır.

Gezegen savunması açısından bu bağ kritiktir: bir asteroidi saptırmak için
ne kadar enerji gerektiği, cismin içinin sıkı bir kaya mı yoksa gevşek bir
moloz yığını mı olduğuna göre kat kat değişir. Dolayısıyla DART'ın asıl
bilimsel kazancı, $\beta$'nın kendisi değil, $\beta$'nın iç yapı hakkında
ne söylediğidir.

\subsection{Literatürdeki boşluk}
Sorun şudur: $\beta$ tek bir sayıdır, iç yapı ise çok boyutludur.
\citet{raducan2024} farklı malzeme kümeleri ve iç yapılar için ileri
benzetimler koşturup gözlemle karşılaştırmış ve düşük yığın yoğunluğunu,
yüzeye yakın ılımlı bir blok kesrini desteklemiştir. Hemen ardından
\citet{senel2025}, \emph{birbirinden farklı} iç yapıların aynı referans
$\beta$ değerini ürettiğini bildirmiştir: bağdaşımı 20--80 Pa aralığında
türdeş bir hedef, çift katmanlı bir iç yapı ve blokların çarpma
noktasında kümelendiği üç katmanlı bir iç yapı.

Bu bir \emph{dejenerasyon} ifadesidir ve projenin çıkış noktasıdır. Eğer
birden fazla iç yapı aynı gözlemi üretiyorsa, "Dimorphos'un bağdaşım
dayanımı şudur" biçiminde bir nokta tahmini vermek, verinin
taşımadığı bir kesinliği iddia etmek olur. Literatürde eksik olan şey bir
tahmin daha değil, şu sorunun nicel cevabıdır:

\begin{sinir}
\textbf{Araştırma sorusu.} $(\alpha_b, Y_0, f)$ üç boyutlu iç yapı
uzayında, DART gözlemi kaç bağımsız \emph{yönü} gerçekten
kısıtlamaktadır --- ve geri kalan yönler için ne söylenemeyeceği nasıl
kayda geçirilir?
\end{sinir}

\subsection{Amaç ve hipotez}
Projenin amacı, bu soruyu cevaplayabilecek bir ileri model ile ön kayıtlı
bir çıkarım yordamını birlikte kurmaktır. Üç iç yapı koordinatı
seçilmiştir:

\begin{itemize}[leftmargin=7mm,itemsep=1pt,topsep=2pt]
  \item $\alpha_b$ --- \textbf{blok gözenekliliği} (distansiyon): cismin
        içindeki iri blokların ne kadar boşluklu olduğu;
  \item $Y_0$ --- \textbf{matris bağdaşım dayanımı}: blokları saran ince
        malzemenin kendini bir arada tutma gücü [Pa];
  \item $f$ --- \textbf{blok hacim kesri}: hacmin ne kadarının iri
        bloklardan oluştuğu.
\end{itemize}

\noindent\textbf{Hipotez.} DART'ın tek skaler gözlemi bu üç koordinatın
\emph{hepsini} ayırt etmeye yetmez; ölçülebilir yön sayısı üçten küçüktür.
Dahası, modelin $\beta$'ya duyarlılığı, çıkarılmak istenen koordinatlar
kadar, araştırmacının kararla sabitlediği model seçimlerine de bağlıdır
--- ve bu seçimlerin etkisi nicel olarak ölçülebilir.

İkinci hipotez, projenin yöntem tarafını belirler: eğer bir model seçimi
gözlemin kendi belirsizliğinden daha büyük bir kayma üretiyorsa, o seçim
sonucun bir parçasıdır ve raporlanmak zorundadır.

% =======================================================================
\section{Yöntem}

\subsection{Hesaplama altyapısı}
Çarpma, düzgünleştirilmiş parçacık hidrodinamiği (SPH;
\citealp{monaghan1992}) ile modellenmiştir. SPH, sürekli ortamı ağ yerine
hareketli parçacıklarla temsil ettiği için, çarpma sonrası ortaya çıkan
büyük şekil değişimlerini ve serbest yüzeyleri doğal biçimde taşır.

Fizik çekirdekleri proje kapsamında yazılmış, GPU'da çalıştırmak için
NVIDIA Warp \citep{macklin2022} kullanılmıştır; Warp yalnızca derleme ve
çekirdek yürütmesini üstlenir, fizik ifadeleri bize aittir. Koşular
TÜBİTAK ULAKBİM TRUBA ulusal hesaplama altyapısındaki NVIDIA H100
GPU'larda yürütülmüştür. Tüm aritmetik FP64'tür. Bileşenler:

\begin{itemize}[leftmargin=7mm,itemsep=1pt,topsep=2pt]
  \item \textbf{Çekirdek fonksiyonu:} Wendland C2 \citep{wendland1995};
  \item \textbf{Durum denklemi:} Tillotson \citep{tillotson1962};
  \item \textbf{Gözeneklilik:} $P$--$\alpha$ gözenek çökmesi
        \citep{carroll1972, jutzi2008};
  \item \textbf{Dayanım:} sürtünmeli Lundborg akma yüzeyi
        \citep{lundborg1968}, içsel sürtünme katsayısı $\mu_f=0{,}6$;
  \item \textbf{Hasar:} Grady--Kipp kırılgan hasar modeli
        \citep{gradykipp1980, benz1995};
  \item \textbf{Öz çekim:} ağaç tabanlı $N$-cisim toplaması.
\end{itemize}

\subsection{Sabit sahne}
Bütün koşular tek bir sabit sahnede yapılır (Şekil~\ref{sek:hedef}).
Hedef, yarı eksenleri $(88{,}5;\ 87{,}0;\ 58{,}0)$ m olan üç eksenli bir
elipsoittir; kalibre edilmiş kütlesi $4{,}2980\times10^{9}$ kg'dır. İç
yapı, matris içine gömülü kürelerle kurulur; yarıçaplar
$r\in[14,56]$ m aralığında ve $\mathrm{d}N/\mathrm{d}r\propto r^{-3}$
dağılımındadır. Çarpıcı üç küreden oluşur: kütlesi 579,4 kg, hızı
6144,9 m\,s$^{-1}$, yüzey normaliyle açısı $17^{\circ}$'dir.

Çözünürlük, çarpma noktası çevresinde iç içe kabuklarla artırılır:
$r=48,\ 24,\ 12,\ 6,\ 3$ m yarıçaplı kabuklarda parçacık aralığı sırasıyla
kabalıktan inceliğe gider. Pilot ve üretim koşularında dış aralık
7,0 m'dir.

\begin{figure}[t]
\centering
\includegraphics[width=\linewidth]{hedef.pdf}
\caption{Sabit pilot ve üretim sahnesi, ekvator düzleminde ($z=0$)
kesilmiş ve \emph{çizilmemiş} --- geometri doğrudan sahne kurucusunun
kendi çıktısından okunmuştur. Alt yarı tel kafes, kesit yüzü ise $f=0{,}25$
değerinde üretim blok üreticisinin yerleştirdiği beş bloğu taşır. Kesit
düzlemi olarak $z=0$ seçilmiştir, çünkü beş bloğun tamamı bu düzlemi
keser; meridyen düzlemi $y=0$ yalnız bir tanesini keser. Teal halkalar,
çarpma noktası çevresindeki iç içe incelme kabuklarıdır.}
\label{sek:hedef}
\end{figure}

\subsection{Benzetimin kapsadığı zaman aralığı}
Bir çarpma olayı, mikrosaniyelik şok cephesinden günler süren fırlatı
yörüngelerine kadar uzanır; hiçbir benzetim bu aralığın tamamını
taşıyamaz. Hangi kısmın taşındığı, hangi kısmın taşınmadığı bu nedenle
bir sonuç kadar önemlidir (Şekil~\ref{sek:kaskad}).

Şok, incelmiş ızgarada taşınır. Kazı akışı, $t_{\mathrm{geçiş}}=1{,}0$ s
olarak sabitlenen \emph{geç evre geçişinde} gelişmesini tamamlar; bu
noktadan sonra hızlı entegrasyon şemasına geçilir. Fırlatı konisi
$t_{\mathrm{son}}=600$ s bitiş noktasında okunur. Bu bitişten sonrası
benzetimin dışındadır: bir parçanın sistemi terk edip etmediği, havada
kalıp kalmadığı ya da yüzeye geri düşüp düşmediği, entegre etmediğimiz
zaman ölçeklerinde belirlenir.

\begin{figure}[t]
\centering
\includegraphics[width=\linewidth]{kaskad.pdf}
\caption{Bir koşunun kapsadığı ve kapsamadığı aralık. Yatay eksen zamanda
logaritmiktir; düşey konum şematiktir. İki kesikli çizgi, veriye
bakılmadan sabitlenmiş iki model zamanıdır. Gri alanın sağı
benzetilmemektedir: bir parçanın üç olası akıbeti orada belirlenir ve bu
rapor hiçbirini sayısal olarak iddia etmez.}
\label{sek:kaskad}
\end{figure}

\subsection{Denetimli deneyler}
Pilot aşamada, tek değişkenin değiştiği beş koşu tamamlanmıştır. Bağdaşım
taraması $(\alpha_b, f)=(1{,}15;\ 0{,}25)$, rastgele tohum 20260906,
parçacık aralığı 7 m, geçiş süresi 1,0 s ve bitiş 600 s sabit tutularak
yapılır. DO1 ve DO2, DY2 referansından yalnız $Y_0$ bakımından ayrılır. DC
koşusu 10 Pa'da seçilmiş bir matris çekme yasasını açar; DN koşusu 10
Pa'da çarpma noktasını 25 m kaydırır. Üretim modeli söz konusu çekme
yasasını \emph{kapalı} tutar; bu bir karardır ve açıkça yazılmıştır.

\subsection{Geçerlilik denetimleri}
Hiçbir koşu, depodaki geçerlilik denetimlerinden geçmeden okunmaz. Dört
yeni dal ve referans DY2 bu denetimleri geçmiştir: enerji sapmaları
$-\%0{,}79$ ile $-\%0{,}80$ arasında, kütle tutarlılık sapmaları
$-\%0{,}05$ ile $-\%0{,}06$ arasında, momentum defteri artıkları ise
$9\times10^{-15}$'ten büyük değildir.

\begin{sinir}
Bu denetimler koşunun \emph{geçerliliğini} destekler;
\emph{çözünürlük yakınsamasını} kanıtlamaz. İkisi farklı şeylerdir ve
raporun hiçbir yerinde yakınsama iddia edilmemektedir.
\end{sinir}

\subsection{Gözlem hedefi}
Benzetim sahnesinin kütlesi, gerçek Dimorphos için bildirilen kütleyle
birebir aynı olmadığından, gözlenen $\beta$ sahne kütlesine göre
düzeltilir. Üretim sahnesi için kullanılan hedef
\begin{equation}
\beta_{\mathrm{göz}}(M_{\mathrm{sahne}})
  = 3{,}5418^{+0{,}188}_{-0{,}247},
  \qquad M_{\mathrm{sahne}} = 4{,}2980\times10^{9}\ \mathrm{kg}
\label{esit:betagoz}
\end{equation}
biçimindedir. Aralığın asimetrisi olabilirlik kurulurken korunur.

\subsection{Üretim tasarımı ve ön kayıt}
Üretim koordinat uzayı $\alpha_b\in[1{,}00;\ 1{,}30]$,
$Y_0\in[1;\ 10^{5}]$ Pa ve $f\in[0{,}05;\ 0{,}50]$'dir. Tasarım, 7,0 m
aralıkta 96 noktalı bir Latin hiperküp ile, bunun içinden önceden
belirlenmiş 12 noktanın 3,5 m aralıkta yinelenmesinden oluşur
(Şekil~\ref{sek:tasarim}). İç içe 12 nokta, deterministik bir kuralla
seçilmiş ve havuzdan tek bir çıktı var olmadan önce bir sınavla
sabitlenmiştir; indisleri $\{11, 20, 33, 44, 45, 57, 58, 62, 74, 78, 87,
92\}$'dir.

\begin{figure}[t]
\centering
% Genişlikler şekillerin DOĞAL boyutuna göre seçildi: büyütme oranı
% ~1,2'yi aşınca şeklin yazıları gövde metninin yanında orantısız duruyor.
\includegraphics[width=0.68\linewidth]{tasarim.pdf}
\caption{Üretim havuzuna bakılmadan kilitlenmiş tasarım,
$(\alpha_b, Y_0)$ yüzünde gösterilmiştir; $f$ aynı hiperküp üzerinde
$[0{,}05;\ 0{,}50]$ aralığında örneklenir. Bu bir \emph{tasarım}
grafiğidir: ölçülmüş hiçbir yanıt ve çıkarılmış hiçbir iç yapı
taşımaz.}
\label{sek:tasarim}
\end{figure}

96 noktalık havuz, 7,0 m aralıkta yaklaşık 480 GPU$\cdot$saat bütçeyle
planlanmıştır. Kaba ve ince çözünürlükler, çok doğruluklu bir vekil model
ile birleştirilir: Kennedy--O'Hagan AR(1) kurgusunda
$y_{\text{orta}}(x) = \rho\, f_{\text{kaba}}(x) + \delta(x)$ yazılır
\citep{kennedy2000}; tasarımın iç içe olması $\delta$'yı tam olarak
hesaplanabilir kılar. Vekil, $\theta$'ya göre dört gruba ayrılmış çapraz
doğrulamayla sınanır ve dışarıda bırakılan artıklar gerektirdiğinde
öngörü varyansı büyütülür \citep{bachoc2013}.

Uygun gözlem vektörü $z$ ve vekil ortalaması $g(\theta)$ için ön kayıtlı
olabilirlik
\begin{equation}
\mathcal{L}(\theta) \propto |S(\theta)|^{-1/2}
  \exp\!\left[-\tfrac12 (g-z)^{\mathsf T} S(\theta)^{-1} (g-z)\right]
\label{esit:olabilirlik}
\end{equation}
biçimindedir; $S(\theta)$ öngörü, gözlem ve kabul edilmiş model
belirsizliğini ölçülen artık korelasyonuyla birleştirir.

\subsubsection*{Karar kapıları}
Projenin yöntemsel omurgası, sonucun ne olacağına veriye bakmadan karar
vermektir. Beş karar kapısının tamamı, ilgili veri daha üretilmemişken
yazılmış, depoya işlenmiş ve tarih damgası almıştır
(Şekil~\ref{sek:kapilar}). Eşikler sonradan ayarlanamaz; iki kapı
yayımlamayı tümüyle durdurur.

\begin{figure}[t]
\centering
\includegraphics[width=\linewidth]{kapilar.pdf}
\caption{Ön kayıtlı karar kapıları. Her kırmızı kutu, ilgili veri var
olmadan \emph{önce} bir protokol belgesinde sabitlenmiş bir koşuldur;
kesikli dallar kendimizi bağladığımız sonuçlardır. Eşikler bize ait
değildir: olasılıksızlık kesmesi $I>3$ \citet{vernon2014}'ün ölçütüdür,
Fisher eşiği ADR-0051'de sabittir ve yanlı bir posterior geçene kadar
olabilirlik paydasını genişletmek yazılı olarak yasaktır. Bir kapı
düşerse, raporlanan şey \emph{daha zayıf} olan sonuçtur.}
\label{sek:kapilar}
\end{figure}

Posterior yayımlanmadan önce geçilmesi gereken denetimler şunlardır:
benzetim tabanlı kalibrasyon (SBC) en az 200 yinelemeyle, gerçek vekil ve
hata modeli kullanılarak yapılır; yanlı bir sıra dağılımı yayımlamayı
engeller \citep{talts2018}. Yerel Fisher bilgisi en az beş parametre
konumunda değerlendirilir; 36'dan büyük öz değerler öğrenilmiş yönleri
sayar ve kayıtlı korelasyon duyarlılıkları üzerindeki \emph{en küçük}
sayı iddiayı belirler. Öncel seçeneği altında $Y_0$ kipinde 0,50 dekattan
büyük bir kayma, tek bir $Y_0$ tahmini verilmesini engeller. Bağımsız
tarih eşleme $I>3$ olan noktaları dışlar \citep{vernon2014}.

% =======================================================================
\section{Bulgular}

\subsection{Tamamlanan koşular}
Tamamlanmış ve geçerliliği denetlenmiş beş koşunun 600 s'deki değerleri
Tablo~\ref{ciz:kosular}'de, dört gözlenebilir üzerindeki yanıt
Şekil~\ref{sek:bulgular}'de verilmiştir. Benzetimin ürettiği fırlatı
kütlesi için $M_{\mathrm{fır}}$ yazılmaktadır; bu büyüklük, gözlenen
toz bulutu kütlesiyle \emph{tanımı gereği} aynı şey değildir.

\begin{table}[t]
\centering
\caption{600 s'de tamamlanan pilot koşular. Belirtilen değişiklik dışındaki
tüm koşullar sabit tutulmuştur. DY2, 10 Pa referans koşusudur.}
\label{ciz:kosular}
\small
\begin{tabular}{@{}llrrrr@{}}
\toprule
Koşu & Değişiklik & $\beta$ & $M_{\mathrm{fır}}$ (kg)
     & Koni ($^\circ$) & $t_{50}$ (s)\\
\midrule
DY2 & $Y_0=10$ Pa           & 3,7480 & $1{,}964\times10^{7}$ & 85,2 & 0,541\\
DO1 & $Y_0=500$ Pa          & 2,6247 & $3{,}649\times10^{6}$ & 77,1 & 0,107\\
DO2 & $Y_0=5000$ Pa         & 1,8738 & $7{,}516\times10^{5}$ & 83,8 & 0,015\\
DC  & Çekme açık            & 1,0807 & $2{,}777\times10^{2}$ & 47,9 & 0,001\\
DN  & Çarpma yeri $+25$ m   & 3,6650 & $2{,}321\times10^{7}$ & 107,7 & 0,845\\
\bottomrule
\end{tabular}
\end{table}

\begin{figure}[t]
\centering
\includegraphics[width=\linewidth]{bulgular.pdf}
\caption{Beş tamamlanmış pilot koşu için dört gözlenebilir atlası. Dolu
daireler ve onları birleştiren doğrular \emph{yalnız} üç bağdaşım
koşusunu (DY2, DO1, DO2) gösterir; doğrular gözü yönlendirmek içindir,
çıkarılmış bir yanıt eğrisi değildir. Kare DC (matris çekmesi açık),
baklava DN (çarpma yeri 25 m kaydırılmış) koşusudur; ikisi de
$Y_0=10$ Pa'da kalır ve bağdaşım taramasının parçası değildir. Koni açısı
ve ortanca fırlatı zamanı $t_{50}$ model tanılarıdır. (a) panelindeki gri
şerit, Denklem~\eqref{esit:betagoz}'in kütleye göre düzeltilmiş DART
$1\sigma$ aralığıdır; posterior aralık değildir. (b) panelinde gözlem
şeridi yoktur, çünkü model fırlatı istatistiği LICIACube toz bulutu
kestirimiyle aynı büyüklük değildir.}
\label{sek:bulgular}
\end{figure}

\subsection{Bağdaşım taraması}
Denenen 10--5000 Pa aralığında $\beta$ 3,7480'den 1,8738'e düşerken,
model fırlatı kütlesi 26,1 kat azalmaktadır. Önceden ilan edilmiş üç
noktalı uyum $\beta-1 = C\,Y_0^{\,p}$ için $p=-0{,}1791$ ve en büyük
bağıl artık 0,1198 bulunmuştur; bu, protokolde sabitlenmiş 0,15'lik uyum
kalitesi eşiğinin altındadır. Denklem~\eqref{esit:betagoz}'deki merkezî
hedefte bu uyum $Y_0\simeq20{,}1$ Pa'ya karşılık gelir
(Şekil~\ref{sek:atlas}).

\begin{figure}[t]
\centering
\includegraphics[width=0.79\linewidth]{atlas.pdf}
\caption{Beş tamamlanmış koşunun bağdaşım--momentum manzarasındaki yeri.
Noktalı eğri, önceden ilan edilmiş üç noktalı kuvvet yasasının çizelgedeki
değerlerden yeniden uydurulmuş hâlidir; düşey çizgi, bu eğrinin kütleye
göre düzeltilmiş DART hedefini kestiği yeri işaretler. Mor şerit
\citet{senel2025}'in bildirdiği türdeş bağdaşım aralığıdır ve
\emph{yalnız yön bulmak için} gösterilmiştir: farklı bir kodla ve farklı
bir içsel sürtünmeyle elde edilmiştir. DC ve DN, DY2 ile aynı
$Y_0=10$ Pa'dadır ve yalnız kendi yaptıkları değişiklik kadar kaymıştır.}
\label{sek:atlas}
\end{figure}

\begin{sinir}
\textbf{20,1 Pa bir ölçüm değildir.} Bu değer; sabit sahneye, sabit
$\alpha_b$ ve $f$'ye, seçilen malzeme yasasına ve uyum biçimine koşullu
bir \emph{aradeğerlemedir}. Dimorphos'un dayanımının posterior tahmini
değildir ve bu raporda öyle kullanılmamaktadır. Tek faydası, eski üretim
alanı $[10^{3}, 10^{7}]$ Pa'nın bu değeri \emph{içermediğini} gösterip
ileriye dönük olarak $[1, 10^{5}]$ Pa alanının seçilmesini
gerekçelendirmesidir.
\end{sinir}

\subsection{Matris çekme dayanımı}
$Y_0=10$ Pa'da, denenen çekme yasasının açılması $\beta$'yı 3,7480'den
1,0807'ye indirir. Fazla momentum $b=\beta-1$ üzerindeki karşıtlık
\begin{equation}
\frac{|b_{\mathrm{DY2}} - b_{\mathrm{DC}}|}{b_{\mathrm{DY2}}}
  = \frac{|2{,}7480 - 0{,}0807|}{2{,}7480} = 0{,}9706
\label{esit:cekme}
\end{equation}
olur. Kayda geçen model fırlatı kütlesi $1{,}964\times10^{7}$ kg'dan
$2{,}777\times10^{2}$ kg'a, yani yaklaşık $7{,}1\times10^{4}$ kat
düşer. Bu, \emph{tek bir bünye modeli seçimine} karşı çok güçlü bir
duyarlılıktır. DC tek bir parametre konumunda tek bir anahtar olduğundan,
gerçek cismin çekme dayanımını belirlemez ve diğer tüm çekme yasalarını
dışlamaz.

\subsection{Çarpma yeri kayması}
Çarpma noktasının 25 m kaydırılması $\beta$'yı 3,7480'den 3,6650'ye
(\%2,2) değiştirirken, model fırlatı kütlesini
$1{,}964\times10^{7}$ kg'dan $2{,}321\times10^{7}$ kg'a (\%18) taşır. Koni
açısı $85{,}2^{\circ}$'den $107{,}7^{\circ}$'ye, ortanca fırlatı zamanı
$t_{50}$ 0,541 s'den 0,845 s'ye çıkar.

Yani \emph{neredeyse sabit bir} $\beta$, bu denetimli çiftte
\emph{belirgin biçimde farklı} bir fırlatı alanıyla bir arada
bulunabilmektedir. Bu, raporun en öğretici gözlemlerinden biridir: $\beta$
tek başına fırlatı alanını belirlemez.

$b$ üzerindeki karşılık gelen çarpma-yeri farkı 0,03019'dur; ön kayıtlı
\emph{ihmal edilebilir} sınırı olan 0,03000'i 0,00019 aşmaktadır.

\begin{sinir}
Sınırı sonradan 0,031'e çekmek sonucu "ihmal edilebilir" yapardı. Kilitli
protokol kuralı gereği eşik \emph{değiştirilmemiş}, sonuç kayıtlı
sınıflandırmayla --- yani sınırı aşmış olarak --- raporlanmıştır. Kuralı
sonuca göre ayarlamamak, bu projenin tek maddelik yöntem sözüdür.
\end{sinir}

\subsection{Ölçülen kaldıraçların sıralaması}
Projenin en bilgilendirici tek grafiği, ölçülmüş bütün kaldıraçları $b$
üzerinde tek bir logaritmik ölçeğe koyan Şekil~\ref{sek:duyarlilik}'tir.

\begin{figure}[t]
\centering
\includegraphics[width=0.94\linewidth]{duyarlilik.pdf}
\caption{$b=\beta-1$ üzerinde ölçülen her kaldıraç, tek logaritmik
ölçekte, büyükten küçüğe. Mavi, çıkarımın öğrenmesi gereken bir
koordinatı; kırmızı, bizim kararla sabitlediğimiz bir model seçimini
gösterir. Yeşil şerit gözlemin $b$ üzerindeki kendi $1\sigma$
genişliğidir; solunda kalan noktalar $\beta$'yı ölçümün ayırt
edebileceğinden daha az oynatır. En büyük iki giriş, tek bir parametre
konumundaki tek denetimli çiftlerden gelir: oradaki duyarlılığı
sınırlarlar, küresel türev değildirler.}
\label{sek:duyarlilik}
\end{figure}

Ön kayıtlı bağıl fark terimleri şunlardır: 0,013 (gerçekleme), 0,004
(uzak alan çözünürlüğü), 0,053 (yakın alan çözünürlüğü), 0,016 (geç zaman
platosu), 0,009 (hedef şekli) ve 0,030 (çarpma yeri). Buna ek olarak
0,134 (çarpıcı geometrisi) ve 0,971 (matris çekmesi), Gauss gürültüsü
değil \emph{seçenek model} duyarlılıklarıdır.

% =======================================================================
\section{Sonuç ve Tartışma}

\subsection{Üç denetimli bulgu}
\textbf{Birincisi,} sabit bir DART sahnesinde, denenen bağdaşım
değerleri boyunca hem momentum artırma katsayısı hem fırlatı kütlesi
azalmaktadır. \textbf{İkincisi,} seçilen matris çekme anahtarı fazla
momentumu \%97,1, model fırlatı istatistiğini ise yaklaşık
$7{,}1\times10^{4}$ kat değiştirmektedir. \textbf{Üçüncüsü,} 25 m'lik bir
çarpma yeri kayması $\beta$'yı yalnız \%2,2 değiştirirken fırlatı
kütlesini, koni açısını ve ortanca fırlatı zamanını daha fazla
değiştirmektedir. Bunlar model duyarlılıklarıdır; yeniden kurulmuş bir iç
yapı değildir.

\subsection{Rahatsız edici sıralama}
Şekil~\ref{sek:duyarlilik}'teki sıralama bu raporun en rahatsız edici
kısmıdır ve bu yüzden gizlenmemiştir. $b$ üzerindeki \emph{en büyük}
ölçülmüş etki, çıkarmaya çalıştığımız koordinatlardan biri değildir:
0,971 ile matris çekme yasasıdır --- yani kararla sabitlediğimiz bir bünye
seçimi. Bu değer, çıkarımın ayırt etmesi beklenen 10--5000 Pa'lık bütün
bağdaşım açıklığını (0,682) aşmaktadır. Çarpıcı geometrisi de 0,134 ile
gözlemin $b$ üzerindeki kendi $1\sigma$ genişliğinin üstündedir. Geri
kalan dört seçim bu genişliğin altında kalır ve eldeki kanıta göre
tartışmaya değmez.

Bunun doğrudan sonucu şudur: $(\alpha_b, Y_0, f)$ üzerindeki bir
posterior, ancak üstünde duran sabit seçimler kadar anlamlıdır. Rapor bu
nedenle bir posterior sunmak yerine o seçimleri kayda geçirmektedir.

\subsection{Literatürle ilişki}
Sahneye koşullu aradeğerlememiz olan 20,1 Pa, \citet{senel2025}'in
türdeş 20--80 Pa aralığının alt kenarına düşmektedir. Bunu bir
\emph{doğrulama} saymıyoruz. İki hesap koda, hedef geometrisine, boyut
sayısına ve --- önemlisi --- içsel sürtünmeye göre farklıdır:
\citet{senel2025} sürtünmeyi $\beta$ üzerinde güçlü bir kaldıraç olarak
bildirmekte ve türdeş aralığını $\mu_f=0{,}5$'te vermektedir; oysa
Tablo~\ref{ciz:kosular}'deki her koşu $\mu_f=0{,}6$ kullanmaktadır.
Farklı sürtünmelerde elde edilmiş iki sayının uyuşması rastlantı olabilir.
Kendi sürtünme duyarlılığımızı henüz ölçmediğimiz için bu karşılaştırma
bir \emph{tutarlılık notundan} ibarettir.

\citet{raducan2024}'un düşük yığın yoğunluğu ve ılımlı blok kesri
sonucuyla sahnemizin kurulumu çelişmemektedir; ancak bu rapor cismin
yoğunluğu üzerine bağımsız bir ölçüm sunmadığından, karşılaştırma yine
yön bulma amaçlıdır.

\subsection{Sınırlar ve iddia etmediklerimiz}
Projenin sınırları yapısaldır ve burada eksiksiz sayılmaktadır.

\begin{enumerate}[leftmargin=7mm,itemsep=2pt,topsep=3pt]
\item \textbf{Üç noktalı uyum sahneye koşulludur.} Tek bir sabit sahnede
      üç noktadan geçirilmiştir ve dışına genellenemez.

\item \textbf{Yakın alan fark terimi koda bağlanmamıştır.} Değer (0,053)
      belgelidir, ama çalıştırılabilir sabiti hâlâ \texttt{nan}'dır ve kod
      bu yapılandırmayı reddetmektedir. Bu bir yuvarlama ayrıntısı
      değildir: $\sigma_{\text{model}}^2$ kareli terimlerin toplamı
      olduğundan, ilan edilen değerler bu tek terimi varyansın
      \%66,4'ü yapar. Terim olmadan $\sigma_{\text{model}}$ 0,0377,
      varken 0,0650 okunur; yani onsuz verilen bir payda karşılaştırmayı
      1,72 kat daraltırdı. Bu nedenle tam bir işletilebilir payda
      \emph{verilmemektedir}.

\item \textbf{Yazılmış ama çağrılmamış modüller vardır.} Depo kusur kütüğü,
      yazılmış, belgelenmiş ve birim sınavlarından geçmiş olduğu hâlde
      üretim yolundan hiç çağrılmayan üç modül kaydetmektedir. Geçen bir
      sınav takımı ile bağlanmış bir üretim yolu aynı özellik değildir ve
      ikincisi için otomatik denetimimiz yoktur.

\item \textbf{Gözlem eşleştirmesi tamamlanmamıştır.} LICIACube
      görüntülerinden kestirilen toz bulutu kütlesi
      $(1{,}6\pm0{,}3)\times10^{7}$ kg, kamera görüş alanındaki kütle için
      makul bir \emph{üst sınır} olarak bildirilmiştir \citep{lolachi2025};
      bizim 600 s'deki model istatistiğimiz tanımı gereği aynı büyüklük
      değildir. Ayrıca mevcut model $\beta$'sı çarpıcı momentumu boyunca
      izdüşürülmekte, yörünge gözlemi ise Dimorphos'un yörünge hız
      doğrultusunu ilgilendirmektedir. İki izdüşüm arasındaki açıyı
      \emph{vermiyoruz}: birincil kaynaktan doğrulamadığımız yörünge
      sabitlerinden çıkar.

\item \textbf{Kaçış eşiği ikiye ayrılmıştır.} Sahnenin kendi kütlesi ve
      etkin yarıçapından cismin kaçış hızı 0,0865 m\,s$^{-1}$ olarak
      hesaplanmaktadır; ancak Didymos \emph{sisteminden} kaçış, ana
      cismin kütlesine ve ayrıklığa bağlıdır ve bunlar birincil kaynaktan
      doğrulanmamıştır. Dimorphos'tan kaçacak kadar hızlı malzeme ikiliye
      bağlı kalıp geri dönebilir.

\item \textbf{Krater tanısı geçersizdir.} Mevcut 600 s bitiş noktasındaki
      krater derinliği, oturmuş bir son kratere fiziksel olarak
      karşılık gelmemektedir.

\item \textbf{Yakınsama kanıtı yoktur.} Daha önceki tek bir çözünürlük
      adımı farklı bir sahnede $\beta-1$ üzerinde \%5,34 ölçmüştür; bu bir
      yakınsama kanıtı değildir.
\end{enumerate}

\begin{sinir}
\textbf{Bu raporda bulunmayanlar.} Üretim posterioru, $(\alpha_b, Y_0, f)$
için herhangi bir nokta tahmini ya da güven aralığı, tanımlanabilir yön
sayısının ölçülmüş değeri ve Hera için sayısal bir öngörü. Bunların
hiçbiri sunulmamaktadır, çünkü Şekil~\ref{sek:kapilar}'deki kapılar henüz
çalıştırılmamıştır. Kapılar veriden önce yazıldığı için, sonuç ne çıkarsa
çıksın kuralı değiştirme imkânımız yoktur --- raporun bu bölümünün
varlık sebebi budur.
\end{sinir}

% =======================================================================
\section{Öneriler}

Aşağıdaki adımlar, raporda sayılan sınırları doğrudan hedefler ve öncelik
sırasıyla verilmiştir.

\begin{enumerate}[leftmargin=7mm,itemsep=2pt,topsep=3pt]
\item \textbf{Sürtünme duyarlılığı ölçülmelidir.} $\mu_f=0{,}5$ ve
      $\mu_f=0{,}6$ arasında denetimli bir çift koşu,
      \citet{senel2025} ile karşılaştırmanın şu anki en zayıf halkasını
      ölçülebilir hâle getirir. Bu, tek bir denetimli çiftle kapatılabilen
      en büyük boşluktur.

\item \textbf{Yakın alan fark terimi koda bağlanmalıdır.} \texttt{nan}
      sabiti, kaynağı belgelenmiş bir değerle değiştirilmeden işletilebilir
      bir olabilirlik paydası verilemez.

\item \textbf{Yörünge sabitleri birincil kaynaktan doğrulanmalıdır.}
      Didymos kütlesi, yörünge yarıçapı ve dönem doğrulanmadan ne izdüşüm
      açısı ne de sistemden kaçış ölçütü hesaplanabilir.

\item \textbf{Çözünürlük, sınırlı bütçeli bir levha yamasıyla
      sınanmalıdır.} Tam cisimde yakınsama taraması bütçeye sığmadığından,
      analitik levha testiyle kapılanmış ayrı bir çözünürlük çalışması
      planlanmıştır; bu çalışmanın yazılı yasağı, sonucunun tam cisme
      aktarıldığını iddia etmemektir.

\item \textbf{Üretim havuzu okunmalı ve kapılar sırayla çalıştırılmalıdır.}
      Havuz tamamlandığında, bu rapordaki hiçbir eşik değiştirilmeden
      Şekil~\ref{sek:kapilar}'deki sıra izlenecek ve sonuç --- posterior,
      daha geniş bir eşleşme bölgesi, ya da hiçbiri --- olduğu gibi
      yazılacaktır.

\item \textbf{Hera karşılaştırması zamana göre planlanmalıdır.} Avrupa
      Uzay Ajansı Hera'nın Didymos sistemine Kasım 2026'da varacağını
      bildirmektedir \citep{esaArrival}. Ancak varış ölçüm, ölçüm de
      yayımlanmış bir sayı demek değildir: DART'ın kendisi, çarpmadan
      yaklaşık iki hafta sonra yörünge dönemi sonucunu, momentum aktarım
      makalelerini ise bundan yaklaşık beş ay sonra vermiştir
      \citep{cheng2023}. Mühürlü bir öngörünün saati her büyüklük için
      ayrı işler ve hepsi varış tarihinden uzundur.
\end{enumerate}

% =======================================================================
\section*{Teşekkür}
Hesaplamalar, TÜBİTAK ULAKBİM Yüksek Başarımlı ve Grid Hesaplama Merkezi
(TRUBA) kaynaklarında yürütülmüştür.

% =======================================================================
\begin{singlespace}
\small
\begin{thebibliography}{99}
\raggedright
\setlength{\itemsep}{1.5pt}

\bibitem[Bachoc(2013)]{bachoc2013}
Bachoc, F. (2013). Cross validation and maximum likelihood estimations of
hyper-parameters of Gaussian processes with model misspecification.
\emph{Computational Statistics \& Data Analysis}, 66, 55--69.

\bibitem[Benz ve Asphaug(1995)]{benz1995}
Benz, W. ve Asphaug, E. (1995). Simulations of brittle solids using smooth
particle hydrodynamics. \emph{Computer Physics Communications}, 87(1--2),
253--265.

\bibitem[Carroll ve Holt(1972)]{carroll1972}
Carroll, M. M. ve Holt, A. C. (1972). Static and dynamic pore-collapse
relations for ductile porous materials. \emph{Journal of Applied Physics},
43(4), 1626--1636.

\bibitem[Cheng ve ark.(2023)]{cheng2023}
Cheng, A. F. ve ark. (2023). Momentum transfer from the DART mission
kinetic impact on asteroid Dimorphos. \emph{Nature}, 616, 457--460.
\url{https://doi.org/10.1038/s41586-023-05878-z}

\bibitem[Daly ve ark.(2023)]{daly2023}
Daly, R. T. ve ark. (2023). Successful kinetic impact into an asteroid for
planetary defence. \emph{Nature}, 616, 443--447.
\url{https://doi.org/10.1038/s41586-023-05810-5}

\bibitem[ESA(2026)]{esaArrival}
Avrupa Uzay Ajansı (2026). \emph{ESA's Hera targets early arrival at
Didymos asteroids}. Erişim: 10 Ekim 2026.
\url{https://www.esa.int/Space_Safety/Hera}

\bibitem[Grady ve Kipp(1980)]{gradykipp1980}
Grady, D. E. ve Kipp, M. E. (1980). Continuum modelling of explosive
fracture in oil shale. \emph{International Journal of Rock Mechanics and
Mining Sciences}, 17(3), 147--157.

\bibitem[Jutzi ve ark.(2008)]{jutzi2008}
Jutzi, M., Benz, W. ve Michel, P. (2008). Numerical simulations of impacts
involving porous bodies I. Implementing sub-resolution porosity in a 3D
SPH hydrocode. \emph{Icarus}, 198(1), 242--255.

\bibitem[Kennedy ve O'Hagan(2000)]{kennedy2000}
Kennedy, M. C. ve O'Hagan, A. (2000). Predicting the output from a complex
computer code when fast approximations are available. \emph{Biometrika},
87(1), 1--13.

\bibitem[Lolachi ve ark.(2025)]{lolachi2025}
Lolachi, R. ve ark. (2025). Ejecta mass estimates from the DART impact
plume inferred from LICIACube images. \emph{The Planetary Science
Journal}, 6, 199. \url{https://doi.org/10.3847/PSJ/adec6b}

\bibitem[Lundborg(1968)]{lundborg1968}
Lundborg, N. (1968). Strength of rock-like materials.
\emph{International Journal of Rock Mechanics and Mining Sciences}, 5(5),
427--454.

\bibitem[Macklin(2022)]{macklin2022}
Macklin, M. (2022). \emph{Warp: A high-performance Python framework for
GPU simulation and graphics}. NVIDIA.

\bibitem[Monaghan(1992)]{monaghan1992}
Monaghan, J. J. (1992). Smoothed particle hydrodynamics. \emph{Annual
Review of Astronomy and Astrophysics}, 30, 543--574.

\bibitem[Raducan ve ark.(2024)]{raducan2024}
Raducan, S. D. ve ark. (2024). Physical properties of asteroid Dimorphos
as derived from the DART impact. \emph{Nature Astronomy}, 8, 445--455.
\url{https://doi.org/10.1038/s41550-024-02200-3}

\bibitem[Senel ve ark.(2025)]{senel2025}
Senel, C. ve ark. (2025). DART-scale impacts into a complex low-strength
target: Insights into asteroid Dimorphos's interior from simulations.
\emph{The Planetary Science Journal}, 6, 157.
\url{https://doi.org/10.3847/PSJ/addf31}

\bibitem[Talts ve ark.(2018)]{talts2018}
Talts, S., Betancourt, M., Simpson, D., Vehtari, A. ve Gelman, A. (2018).
\emph{Validating Bayesian inference with simulation-based calibration}.
arXiv:1804.06788.

\bibitem[Tillotson(1962)]{tillotson1962}
Tillotson, J. H. (1962). \emph{Metallic equations of state for
hypervelocity impact} (Rapor No. GA-3216). General Atomic.

\bibitem[Vernon ve ark.(2014)]{vernon2014}
Vernon, I., Goldstein, M. ve Bower, R. G. (2014). Galaxy formation:
Bayesian history matching for the observable universe. \emph{Statistical
Science}, 29(1), 81--90.

\bibitem[Wendland(1995)]{wendland1995}
Wendland, H. (1995). Piecewise polynomial, positive definite and compactly
supported radial functions of minimal degree. \emph{Advances in
Computational Mathematics}, 4(1), 389--396.

\end{thebibliography}
\end{singlespace}

\end{document}
